\documentclass[11pt]{article}
\usepackage[T1]{fontenc}
\usepackage{lmodern,amsmath,amssymb,amsthm,mathtools}
\usepackage[margin=1.05in]{geometry}
\usepackage[hidelinks]{hyperref}
\newtheorem{theorem}{Theorem}[section]
\newtheorem{lemma}[theorem]{Lemma}
\newtheorem{proposition}[theorem]{Proposition}
\newtheorem{corollary}[theorem]{Corollary}
\newcommand{\R}{\mathbb R}\newcommand{\E}{\mathbb E}
\newcommand{\Hh}{\mathcal H}\newcommand{\tr}{\operatorname{tr}}
\newcommand{\code}[1]{{\small\nolinkurl{#1}}}
\title{Diagonal Gaussian forcing:\\trace thresholds, completed divergence and recovery geometry}
\author{Project Stardust}
\date{Reviewed formalization companion, October 29, 2026}
\begin{document}
\maketitle
\begin{abstract}
Countably many independent exponentially damped Gaussian modes define
an actual square-integrable Hilbert-valued field at positive time exactly
when the squared forcing amplitudes divided by the damping rates are
summable. We prove this criterion, the exact covariance and cutoff error,
and mean-square continuity up to time zero. The field is a closed linear
Gaussian field in the original source space; its full Gaussian divergence
exists, lies in second chaos, and has norm at most $\sqrt2$ times the
field norm. The proof reconstructs finite Gaussian integration by parts
and Hilbert--Schmidt completion, including the symmetric/antisymmetric
kernel distinction. In contrast, observation of all Brownian modes at
any positive time retains an infinite first-chaos isometry, preventing
compact divergence recovery. Finite field energy and compact recovery
are therefore separate questions. The trace and Gaussian mechanisms are
classical; the paper consolidates frozen verified constructions and
explicit counterexamples without claiming new stochastic PDE theory.
\end{abstract}

\section{The actual modes and the trace criterion}
Let $\Hh$ be a real separable infinite-dimensional Hilbert space with
an isonormal Gaussian map $X:\Hh\to L^2(P)$:
$\E X(h)=0$, $\E[X(h)X(k)]=\langle h,k\rangle$, and finite families
are jointly Gaussian. All scalar and field $L^2$ spaces below are over
the Gaussian sigma-algebra generated by $X(\Hh)$, modulo null sets;
no additional non-Gaussian random factor is included. Embed mutually
orthogonal copies of
$L^2(\R_+)$ isometrically into $\Hh$, with embeddings $\iota_n$,
$n\geq0$. Thus $B_n(t)=X(\iota_n\mathbf1_{(0,t]})$ are independent
Brownian coordinates at the level of their joint Gaussian laws.
Continuous versions give the usual Brownian realization if desired;
only the stated $L^2$ coordinates are used below.

Fix deterministic rates $\lambda_n\geq1$, amplitudes $c_n\in\R$,
and an orthonormal sequence $(e_n)$ in the same source coefficient
space $\Hh$. Define
\begin{equation}\label{eq:profiles}
 k_{n,t}(s)=\mathbf1_{(0,t]}(s)e^{-\lambda_n(t-s)},\quad
 h_{n,t}=\iota_nk_{n,t},\quad Y_n(t)=X(h_{n,t}).
\end{equation}
The deterministic-integrand Brownian notation
$Y_n(t)=\int_0^te^{-\lambda_n(t-s)}\,dB_n(s)$ is equivalent to this
isonormal definition. No adaptedness of the coefficient vectors $e_n$
relative to the noise is needed. They are deterministic output
directions; they need not be orthogonal to all input directions.

For $N\geq1$ put
\[
 F_N(t)=\sum_{n<N}c_n e_nY_n(t)\in L^2(P;\Hh).
\]
This is a Hilbert-valued random field, not a scalar sum and not yet a
claimed solution of an arbitrary SPDE.

\begin{theorem}\label{thm:trace}
For each $t>0$, the following are equivalent:
\begin{enumerate}
\item $F_N(t)$ converges in $L^2(P;\Hh)$;
\item $\displaystyle\sum_{n\geq0}\frac{c_n^2}{\lambda_n}<\infty$.
\end{enumerate}
When these hold, write the limit as $F(t)$. Then
\begin{align}
 \|F(t)\|_2^2
 &=\sum_{n\geq0}c_n^2\frac{1-e^{-2\lambda_nt}}{2\lambda_n},
 \label{eq:energy}\\
 \|F(t)-F_N(t)\|_2^2
 &=\sum_{n\geq N}c_n^2\frac{1-e^{-2\lambda_nt}}{2\lambda_n}
 \leq\frac12\sum_{n\geq N}\frac{c_n^2}{\lambda_n}.
 \label{eq:tail}
\end{align}
Setting $F(0)=0$, the curve $t\mapsto F(t)$ is continuous in
$L^2(P;\Hh)$ on $[0,\infty)$.
\end{theorem}

The reciprocal-rate condition is a diagonal instance of the classical
stochastic-convolution Hilbert--Schmidt criterion; see Hairer's
primary lecture notes~\cite{Hairer}, Sections 4.4 and 6. The theorem
here fixes the actual source construction and distinguishes its
positive-time quantifier from the trivial zero-time field.

\section{Finite-mode identities and completion}
\begin{lemma}\label{lem:profile}
For $\lambda>0$, let $k_{\lambda,t}(r)=\mathbf1_{(0,t]}(r)
e^{-\lambda(t-r)}$ and
$V_\lambda(t)=(1-e^{-2\lambda t})/(2\lambda)$. Then
\begin{align}
 \|k_{\lambda,t}\|_{L^2}^2&=V_\lambda(t),\label{eq:V}\\
 \langle k_{\lambda,t},k_{\lambda,s}\rangle
 &=e^{-\lambda|t-s|}V_\lambda(\min(s,t)),\label{eq:cov}\\
 \|k_{\lambda,t}-k_{\lambda,s}\|_{L^2}^2
 &=V_\lambda(t)+V_\lambda(s)
       -2e^{-\lambda|t-s|}V_\lambda(\min(s,t)).\label{eq:increment}
\end{align}
For $\lambda\geq1$,
\begin{equation}\label{eq:bounds}
 \frac{1-e^{-2t}}{2\lambda}\leq V_\lambda(t)\leq\frac1{2\lambda},
 \qquad \|k_{\lambda,t}-k_{\lambda,s}\|^2\leq\frac2\lambda.
\end{equation}
\end{lemma}
\begin{proof}
Integrating $e^{-2\lambda(t-r)}$ over $0<r\leq t$ proves
\eqref{eq:V}. If $s\leq t$, factor
$e^{-\lambda(t-s)}$ from the overlap integral on $(0,s]$ and use
\eqref{eq:V}; symmetry gives \eqref{eq:cov}. Expansion of a squared
norm proves \eqref{eq:increment}. The first bounds follow from
$0\leq e^{-2\lambda t}\leq e^{-2t}\leq1$. Finally
\[
 \|k_{\lambda,t}-k_{\lambda,s}\|^2
 \leq2\|k_{\lambda,t}\|^2+2\|k_{\lambda,s}\|^2\leq2/\lambda.
\]
\end{proof}

\begin{proof}[Proof of Theorem~\ref{thm:trace}]
Orthonormality of the deterministic output vectors gives, for any
finite index set $A$,
\begin{equation}\label{eq:finiteenergy}
 \left\|\sum_{n\in A}c_ne_nY_n(t)\right\|_2^2
 =\sum_{n\in A}c_n^2\E Y_n(t)^2
 =\sum_{n\in A}c_n^2V_{\lambda_n}(t).
\end{equation}
In particular, disjoint finite index sets give orthogonal fields.
Thus the partial sums are Cauchy exactly when the nonnegative energy
series converges. At fixed $t>0$, the factor $1-e^{-2t}$ is positive,
so \eqref{eq:bounds} makes that convergence equivalent to the
reciprocal-rate series. Completeness gives $F(t)$.
Passing finite orthogonal sums to their limits gives \eqref{eq:energy}
and \eqref{eq:tail}, including the stated upper error bound.

For continuity fix $s\geq0$. Each right side in
\eqref{eq:increment} is continuous in $t$ and vanishes at $t=s$.
Consequently $k_{\lambda_n,t}\to k_{\lambda_n,s}$ in $L^2$ for each
$n$. Orthogonality gives
\begin{equation}\label{eq:fieldincrement}
 \|F(t)-F(s)\|_2^2
 =\sum_{n\geq0}c_n^2\|k_{\lambda_n,t}-k_{\lambda_n,s}\|_{L^2}^2.
\end{equation}
For example, this equality follows first for $F_N$, then by its
convergence at $t$ and $s$ and the dominated summability in
\eqref{eq:bounds}. The summable bound $2c_n^2/\lambda_n$ permits
dominated convergence of the series, so the squared increment norm
tends to zero. This also proves continuity at $s=0$.
\end{proof}

When $c_n=1$, the rates $(n+1)^2$ admit the field at all times, while
$\lambda_n=n+1$ do not admit the $L^2$ series at any positive time.
Indeed the reciprocal square series converges, for example by
comparison with the telescoping series $1/(n(n+1))$ for $n\geq1$;
the harmonic series diverges by grouping indices in dyadic blocks.
For an undamped unit Brownian series the cutoff energy is $Nt$.
Zero damping lies outside the hypothesis $\lambda_n\geq1$ and is
this separate elementary example, not a substitution into
\eqref{eq:energy} with a zero denominator.

These statements assert $L^2$ convergence and mean-square continuity.
They do not assert a samplewise uniform-in-time limit, a path-space
H\"older exponent, or a continuum nonlinear PDE solution.

\section{The original Gaussian divergence on linear fields}
Let $D$ be the original closed Gaussian gradient, obtained by closing
the polynomial cylindrical gradient in $L^2(P)\to L^2(P;\Hh)$.
Its maximal Hilbert adjoint is the Gaussian divergence $\delta=D^*$.
The relevant conventions agree with~\cite{SV}. Write
$\mathcal L$ for the closed linear span in $L^2(P;\Hh)$ of fields
$aX(b)$ with $a,b\in\Hh$. This is the space of linear Gaussian fields,
not the full random field space.

\begin{lemma}\label{lem:rankone}
For $a,b\in\Hh$,
\begin{equation}\label{eq:rankone}
 aX(b)\in D(\delta),\qquad
 \delta(aX(b))=X(a)X(b)-\langle a,b\rangle.
\end{equation}
For a finite-rank deterministic operator
$K=\sum_{j=1}^m a_j\otimes b_j$, where
$(a\otimes b)h=a\langle b,h\rangle$, put
$F_K=\sum_j a_jX(b_j)$. Then
\begin{align}
 \|F_K\|_2^2&=\|K\|_{\mathrm{HS}}^2,\label{eq:HS}\\
 \|\delta F_K\|_2^2
 &=\|K\|_{\mathrm{HS}}^2+\langle K,K^*\rangle_{\mathrm{HS}}
 =2\|\operatorname{sym}K\|_{\mathrm{HS}}^2
 \leq2\|F_K\|_2^2.\label{eq:quadratic}
\end{align}
Here $\operatorname{sym}K=(K+K^*)/2$. The divergence is in second
Gaussian chaos.
\end{lemma}
\begin{proof}
For a polynomial cylindrical test $f$, Gaussian integration by parts
applied to the product $fX(b)$ gives
\[
 \E[X(a)fX(b)]=\E[\langle Df,a\rangle X(b)]
                       +\langle a,b\rangle\E f.
\]
This proves \eqref{eq:rankone} on the core. Passing graph approximants
of $f$ through the pairing extends it to every $f\in D(D)$; the
right side of \eqref{eq:rankone} is in $L^2$ by Gaussian moments.
It therefore meets the defining maximal-adjoint bound.

The isonormal covariance gives
$\E\langle F_K,F_L\rangle=\langle K,L\rangle_{\mathrm{HS}}$,
proving \eqref{eq:HS} and representation independence. For the
divergence, the centered Gaussian fourth-moment identity gives
\begin{align*}
 &\E[(X(a)X(b)-\langle a,b\rangle)
       (X(c)X(d)-\langle c,d\rangle)]\\
 &\hspace{2em}=\langle a,c\rangle\langle b,d\rangle
                       +\langle a,d\rangle\langle b,c\rangle.
\end{align*}
That identity follows by differentiating the finite joint Gaussian
exponential moment formula four times at zero. Summing proves the
first equality of \eqref{eq:quadratic}. Orthogonality of the symmetric
and skew-adjoint parts in Hilbert--Schmidt inner product proves the
second equality and the upper bound. The centered quadratic products
are degree-two Hermite polynomials or their polarizations, hence lie
in second chaos.
\end{proof}

\begin{theorem}\label{thm:complete}
Every field $F\in\mathcal L$ belongs to the original domain $D(\delta)$.
The restricted divergence is bounded, with
\begin{equation}\label{eq:sqrt2}
 \|\delta F\|_2\leq\sqrt2\|F\|_2.
\end{equation}
It maps into second chaos. For $F=F_K$ with $K$ Hilbert--Schmidt,
its kernel consists exactly of skew-adjoint $K$. Its canonical
gradient recovery is $F_{\operatorname{sym}K}$ and
\begin{equation}\label{eq:recovery}
 R(\delta F_K)=F_{\operatorname{sym}K},\qquad
 \|\delta F_K\|_2^2=2\|F_{\operatorname{sym}K}\|_2^2.
\end{equation}
\end{theorem}
\begin{proof}
The finite-rank isometry \eqref{eq:HS} extends by Hilbert--Schmidt
completion onto $\mathcal L$. For finite-rank approximants $K_j\to K$,
Lemma~\ref{lem:rankone} shows that $\delta F_{K_j}$ is Cauchy, with
the bound \eqref{eq:sqrt2}. Closedness of the original adjoint
therefore puts $(F_K,\lim_j\delta F_{K_j})$ in its original graph.
Second chaos is closed, so the limit remains in it. The exact norm
formula \eqref{eq:quadratic} passes to the limit, proving the kernel
and energy assertions.

On centered Gaussian chaos define the canonical recovery as
$R=DN^{-1}$, where the number operator $N$ has eigenvalue $q$ on
chaos $q\geq1$, and set $R1=0$. On second chaos it is $D/2$.
Finite differentiation gives
\[
 D\delta F_K=F_{K+K^*}=2F_{\operatorname{sym}K}
\]
for finite-rank $K$. Passing to the original closed gradient graph
proves the identity for Hilbert--Schmidt $K$, since both sides
converge strongly. This gives \eqref{eq:recovery}.
\end{proof}

For clarity, the recovery used here is bounded on full scalar $L^2$:
if $f=\sum_{q\geq0}f_q$ is its orthogonal chaos decomposition, the
finite Gaussian identity
$\|Df_q\|_2^2=q\|f_q\|_2^2$ gives
\[
 \|Rf\|_2^2=\sum_{q\geq1}\frac{\|f_q\|_2^2}{q}
 \leq\|f\|_2^2.
\]
The finite identity follows from the one-coordinate Hermite derivative
and orthogonality, then tensorization. Different differentiated chaos
orders are orthogonal. Thus the series defines $R$ by norm completion.
In the frozen library it is constructed as the minimal-norm original
divergence recovery; its chaos identity identifies it with this $DN^{-1}$.
Only its proved action in \eqref{eq:recovery} and on first chaos is
needed in this paper.

The constant $\sqrt2$ is attained for a unit $a=b$ in
\eqref{eq:rankone}. Antisymmetric finite tensors instead give zero
divergence with a nonzero field. Thus no isometry on all of
$\mathcal L$ is asserted, and scalar divergence does not recover the
antisymmetric component of a field.

\section{Divergence and limits of the admitted forcing}
Under the trace criterion of Theorem~\ref{thm:trace}, define
\[
 K_t=\sum_{n\geq0}c_ne_n\otimes h_{n,t}.
\]
Orthonormality of $e_n$ makes this an actual Hilbert--Schmidt sum:
$\|K_t\|_{\mathrm{HS}}^2=\sum_n c_n^2\|h_{n,t}\|^2$.
It follows that $F(t)=F_{K_t}\in\mathcal L$.

\begin{corollary}\label{cor:forcing}
For every admitted forcing profile, $F(t)\in D(\delta)$ at every
$t\geq0$, and
\[
 \delta F_N(t)\longrightarrow\delta F(t)\quad\hbox{in }L^2(P),
 \qquad \|\delta F(t)-\delta F_N(t)\|_2
          \leq\sqrt2\|F(t)-F_N(t)\|_2.
\]
The curve $t\mapsto\delta F(t)$ is mean-square continuous and has the
exact symmetric recovery \eqref{eq:recovery} at each time.
\end{corollary}
\begin{proof}
Apply Theorem~\ref{thm:complete} to $F(t)$ and to the differences of
finite or limiting fields in $\mathcal L$. The norm estimate proves
the convergence. Applying the same estimate to $F(t)-F(s)$ and using
Theorem~\ref{thm:trace} proves continuity. Recovery follows from the
identity for $K_t$.
\end{proof}

This divergence is the original Gaussian source adjoint, not an It\^o
integral in the physical time of $F$. In particular the field's
deterministic output coefficient vectors may overlap its source
directions, and the centered quadratic correction in
\eqref{eq:rankone} must be retained. The trace criterion controls the
specific field; it does not make every observation or recovery map
compact.

\section{Infinitely many observed modes and noncompact recovery}
Let the complete modal past at time $t$ be
\[
 \mathcal A_t=\sigma(B_n(s):n\geq0,\ 0\leq s\leq t),\qquad
 Q_t=\E[\,\cdot\mid\mathcal A_t].
\]
These are the actual Gaussian observation spaces of all modes,
independent of whether a particular amplitude/rate profile is
trace-admitted.

\begin{theorem}\label{thm:noncompact}
The bounded operator $RQ_t:L^2(P)\to L^2(P;\Hh)$ is compact if and
only if $t=0$. At every positive $t$, even recovery restricted to the
observed first chaos contains an infinite-dimensional isometry.
\end{theorem}
\begin{proof}
For $t>0$ set $u_n=t^{-1/2}\iota_n\mathbf1_{(0,t]}$. These form an
orthonormal sequence in $\Hh$, and $X(u_n)$ form an orthonormal
sequence in $L^2(P)$. Each is observed at time $t$, so
$Q_tX(u_n)=X(u_n)$. On first chaos the number operator is one and
$DX(u_n)=u_n$, a deterministic coefficient vector. Therefore
\[
 RQ_tX(u_n)=u_n\quad\hbox{as a constant field in }L^2(P;\Hh).
\]
These outputs are pairwise at distance $\sqrt2$. A compact operator
would take this bounded input sequence to a sequence with a norm
convergent subsequence, a contradiction. Taking closed linear spans
also proves the asserted isometry. At $t=0$ all modal observations
are zero; $Q_0$ projects onto constants, which $R$ kills. Thus
$RQ_0=0$ is compact.
\end{proof}

The obstruction does not depend on the chosen canonical right inverse.
Indeed, if a bounded linear map $L$ satisfies
$\delta L(X(u))=X(u)$ for every $u$ in an infinite observed coefficient
subspace, adjointness against $X(v)$ gives
\[
 \langle\E L(X(u)),v\rangle
 =\E\langle L(X(u)),DX(v)\rangle=\langle u,v\rangle.
\]
Thus expectation, followed by projection to that subspace, is a
bounded left inverse to $u\mapsto L(X(u))$. Compactness of $L$ would
make the identity on an infinite-dimensional Hilbert space compact,
contradicting its orthonormal sequence. This is the elementary
left-inverse mechanism behind the frozen general recovery obstruction.

At zero damping, unit cutoffs have divergent energy; under quadratic
damping, the particular field has finite energy and a valid divergence.
Nevertheless the underlying observation family still contains every
$B_n(t)$ and retains Theorem~\ref{thm:noncompact}. The two conclusions
concern different operators and do not contradict one another.

\section{Formal correspondence and prior theory}
\begingroup\raggedright
The source authority is frozen verified commit
\code{bab688bb6140cd6316c80680095a914ea0d346db}.
The countable isonormal modes and actual orthogonal series are in
\code{SPDECountableModes}; exponential profiles and reciprocal-rate
criteria are in \code{SPDEExponentialForcing}. The actual Hilbert-valued
field and divergence application are in
\code{SPDEHilbertExponentialForcing}; time continuity is
\code{SPDEExponentialTimeContinuity}. The completed original
divergence, second-chaos and projection-energy results are
\code{GaussianLinearFieldDivergenceCompletion}. The full observed
mode obstruction is \code{SPDEObservationRecoveryCompactness}.

The theorem map distinguishes exact declarations from elementary human
reorganizations: Hilbert--Schmidt tensors present the formal closed
linear-field space, the finite fourth-moment proof reconstructs its
quadratic completion, and cutoff errors follow from the exact
orthogonal energy identities. The maximal Gaussian adjoint and
coefficient-space geometry remain those of the frozen construction.
No coefficient alignment is silently added to the application.

Hairer's primary notes establish the classical Hilbert--Schmidt
stochastic-integral and stochastic-convolution setting. The scalar
reciprocal-rate threshold, finite-chaos divergence and infinite-mode
compactness obstruction are presented here as a consolidated
formalization companion. Gaussian derivative/adjoint conventions are
also compared with the primary treatment~\cite{SV}. No new stochastic
PDE existence theorem, optimal sample-path regularity, or priority
claim is made. A samplewise space-time realization, nonlinear equation
or stronger convergence topology requires its own hypotheses and
proof. Full baseline build and axiom evidence was hash-checked rather
than repeated at the expense of active research resources.
\par\endgroup

\begin{thebibliography}{9}
\bibitem{Hairer} M. Hairer,
\emph{An introduction to stochastic PDEs},
\href{https://arxiv.org/abs/0907.4178}{arXiv:0907.4178},
version 2 (2023), Sections 4.4 and 6.
\bibitem{SV} T. Sottinen and L. Viitasaari,
\emph{Stochastic analysis of Gaussian processes via Fredholm representation},
Int. J. Stoch. Anal. (2016), article 8694365.
\href{https://doi.org/10.1155/2016/8694365}{doi:10.1155/2016/8694365}.
\end{thebibliography}
\end{document}
